
\documentclass[12pt, a4paper]{report}

% --- UNIVERSAL PREAMBLE BLOCK ---
\usepackage[a4paper, top=3cm, bottom=3cm, left=4cm, right=3cm]{geometry}
\usepackage[utf8]{inputenc}
\usepackage{lmodern}
\usepackage[indonesian]{babel}
\usepackage{setspace}
\usepackage{titlesec}
\usepackage{tabularx}
\usepackage{graphicx}
\usepackage[normalem]{ulem}
\usepackage{amsmath}


% Pengaturan format Judul Bab
\titleformat{\chapter}[display]
  {\normalfont\large\bfseries\centering}{\MakeUppercase{\chaptertitlename} \thechapter}{5pt}{\large}
\titlespacing*{\chapter}{0pt}{-20pt}{20pt}

% Mendefinisikan kolom C agar teks di-wrap tapi tetap ditengah (centering)
% Kolom X standar itu left-aligned (justify)
\newcolumntype{C}{>{\centering\arraybackslash}X}
\newcolumntype{L}{>{\raggedright\arraybackslash}X}

% --------------------------------------------------------
% DEKLARASI MAKRO UNTUK FORMAT HALAMAN AWAL
% --------------------------------------------------------
\makeatletter
\newcommand{\judul}[1]{\gdef\@judul{#1}}
\newcommand{\jenisTA}[1]{\gdef\@jenisTA{#1}}
\newcommand{\thselesai}[1]{\gdef\@thselesai{#1}}
\newcommand{\tglpenulisan}[1]{\gdef\@tglpenulisan{#1}}
\newcommand{\tglujian}[1]{\gdef\@tglujian{#1}}
\newcommand{\namaA}[1]{\gdef\@namaA{#1}}
\newcommand{\nimA}[1]{\gdef\@nimA{#1}}
\newcommand{\namaB}[1]{\gdef\@namaB{#1}}
\newcommand{\nimB}[1]{\gdef\@nimB{#1}}
\newcommand{\namaC}[1]{\gdef\@namaC{#1}}
\newcommand{\nimC}[1]{\gdef\@nimC{#1}}
\newcommand{\namaD}[1]{\gdef\@namaD{#1}}
\newcommand{\nimD}[1]{\gdef\@nimD{#1}}
\newcommand{\koordinator}[1]{\gdef\@koordinator{#1}}
\newcommand{\nipkoordinator}[1]{\gdef\@nipkoordinator{#1}}
\newcommand{\praktisiA}[1]{\gdef\@praktisiA{#1}}
\newcommand{\praktisiB}[1]{\gdef\@praktisiB{#1}}
\newcommand{\kajur}[1]{\gdef\@kajur{#1}}
\newcommand{\nipkajur}[1]{\gdef\@nipkajur{#1}}

\def\halamansampul{%
  \thispagestyle{empty}%
  \begin{center}
    {\normalfont\bfseries \@jenisTA}\par\nobreak\vspace{1.0cm}
    {\normalfont\large\bfseries\MakeUppercase{\@judul}}\par\nobreak\vspace{1.5cm}
    {\normalfont 
        untuk memenuhi sebagian persyaratan\\
        untuk memperoleh derajat Sarjana S1\\
        Program Studi Teknik Fisika \space} 
    \vfill
    \includegraphics[width=5cm]{logoUGM.png}
    \vfill\vspace{1.0cm}
    {\normalfont Disusun oleh Kelompok:}\\[0.5cm]
    \begin{tabular}{ll}
      1. \@namaA & (\@nimA)\\
      2. \@namaB & (\@nimB)\\
      3. \@namaC & (\@nimC)\\
      4. \@namaD & (\@nimD)
    \end{tabular}\\
    \vspace{1.5cm}
    \begin{onehalfspace}
    {\normalfont\bfseries DEPARTEMEN TEKNIK NUKLIR DAN TEKNIK FISIKA\\FAKULTAS TEKNIK\\UNIVERSITAS GADJAH MADA\\YOGYAKARTA\\\@thselesai}
    \end{onehalfspace}
  \end{center}\clearpage
}

\def\halamanpengesahan{%
  \clearpage\thispagestyle{empty}%
  \addcontentsline{toc}{chapter}{LEMBAR PENGESAHAN}
  \begin{onehalfspace}
    \begin{center}
      {\normalfont\large\bfseries LEMBAR PENGESAHAN}\par\nobreak\vspace{0.5cm}
      {\normalfont \@jenisTA}\par\nobreak\vspace{0.5cm}
      {\normalfont\bfseries\MakeUppercase{\@judul}}\par\nobreak\vspace{0.5cm}
    \end{center}
    \noindent Disusun oleh Kelompok:\\[0.2cm]
    \begin{tabularx}{\textwidth}{l L l}
      1. & \@namaA & (\@nimA) \\
      2. & \@namaB & (\@nimB) \\
      3. & \@namaC & (\@nimC) \\
      4. & \@namaD & (\@nimD) \\
    \end{tabularx}
    \vspace{0.5cm}
    \begin{center} Telah diperiksa dan disetujui pada tanggal \@tglujian \end{center}
    \vspace{0.5cm}
    
    \begin{center}
      \textbf{Tim Penilai:}\\[0.5cm]
       \begin{tabularx}{\textwidth}{l l l}
      1. & Dosen Koordinator: & \@koordinator \\
      2. & Praktisi Industri: & \@praktisiA \\
      3. & Praktisi Industri: & \@praktisiB  \\
    \end{tabularx}
    \end{center}
    
    \vspace{1.0cm}
    \begin{center}
      Mengetahui, \\
      Ketua Departemen Teknik Nuklir dan Teknik Fisika \\ Fakultas Teknik UGM \\[2.2cm]
      \uline{\@kajur} \\ NIP. \@nipkajur
    \end{center}
  \end{onehalfspace}\clearpage
}

\def\paktaintegritas{%
  \clearpage\thispagestyle{empty}%
  \addcontentsline{toc}{chapter}{PAKTA INTEGRITAS}
  \begin{onehalfspace}
  \begin{center} {\normalfont\large\bfseries PAKTA INTEGRITAS}\par\nobreak \end{center}
  \vspace{1cm}
  \noindent Kami yang bertanda tangan di bawah ini (Tim Penyusun PSTF 1):
  \par\nobreak\vspace{0.5cm}
  
  \noindent
  \begin{tabularx}{\textwidth}{l L l}
    1. & \@namaA & (\@nimA) \\
    2. & \@namaB & (\@nimB) \\
    3. & \@namaC & (\@nimC) \\
    4. & \@namaD & (\@nimD) \\
  \end{tabularx}
  
  \vspace{0.5cm}
  \noindent menyatakan dengan sesungguhnya bahwa Laporan Perancangan Sistem Teknik Fisika 1 (PSTF 1) ini adalah murni hasil karya rancangan tim kami sendiri dan tidak mengandung unsur plagiasi. Seluruh data, referensi dokumen lelang (TOR), kode dan standar, dan parameter desain yang digunakan telah disitasi dengan benar. Apabila di kemudian hari terbukti ada pelanggaran akademik, kami bersedia menerima sanksi yang berlaku.
  
  \vspace{0.5cm}
  \noindent\hfill Yogyakarta, \@tglpenulisan
  
  \vspace{0.5cm}
  \noindent
  \begin{tabularx}{\textwidth}{C C}
     Anggota 1 & Anggota 2 \\[1.5cm]
     \uline{\@namaA} & \uline{\@namaB} \\ \@nimA & \@nimB
  \end{tabularx}
  
  \vspace{0.5cm}
  \noindent
  \begin{tabularx}{\textwidth}{C C}
     Anggota 3 & Anggota 4 \\[1.5cm]
     \uline{\@namaC} & \uline{\@namaD} \\ \@nimC & \@nimD
  \end{tabularx}
  
  \vspace{1cm}
  \begin{center} 
    \fbox{\begin{minipage}{0.25\textwidth} \centering \vspace{0.5cm} \textit{Materai ASLI (e-Meterai)}\par \par \textbf{Rp. 10.000} \vspace{0.5cm} \end{minipage}} 
  \end{center}
  \end{onehalfspace}\clearpage
}

\def\ringkasaneksekutif{%
    \clearpage\thispagestyle{empty}\addcontentsline{toc}{chapter}{RINGKASAN EKSEKUTIF}
    \begin{center} {\large\bfseries\MakeUppercase{\@judul}}\par\nobreak\vskip 6\p@ {\normalsize oleh kelompok:}\par\nobreak\vskip 6\p@
        \begin{singlespace} {\normalsize \@namaA \ (\@nimA) \\ \@namaB \ (\@nimB) \\ \@namaC \ (\@nimC) \\ \@namaD \ (\@nimD)}\vskip 10\p@ \end{singlespace}
    \end{center} \vspace{1cm} {\parindent \z@ \centering \normalfont {\normalsize\bfseries RINGKASAN EKSEKUTIF}\par\nobreak } \vspace{0.5cm} \begin{onehalfspace}
}
\def\endringkasaneksekutif{\end{onehalfspace}\clearpage}
\makeatother

% ==========================================
% DATA IDENTITAS
% ==========================================
\judul{PERANCANGAN SISTEM INSTRUMENTASI DAN KONTROL PROSES ZXY PADA PLANT XYZ}
\jenisTA{LAPORAN PERANCANGAN SISTEM TEKNIK FISIKA 1}

\namaA{Nama Anggota 1} \nimA{19/123451/TK/00001}
\namaB{Nama Anggota 2} \nimB{19/123452/TK/00002}
\namaC{Nama Anggota 3} \nimC{19/123453/TK/00003}
\namaD{Nama Anggota 4} \nimD{19/123454/TK/00004}

\thselesai{2026}
\tglpenulisan{10 Mei 2026}
\tglujian{20 Mei 2026}

\koordinator{Nama Dosen Koordinator, S.T., M.T.} \nipkoordinator{197X0000000000}
\praktisiA{Nama Praktisi Industri Satu, S.T.} \praktisiB{Nama Praktisi Industri Dua, S.T., M.T.}
\kajur{Prof. Dr. Ir. Faridah, S.T., M.Sc., IPU. ASEAN Eng} \nipkajur{197602142002122001}

% ==========================================
% MULAI ISI DOKUMEN
% ==========================================
\begin{document}
\pagenumbering{roman}
\halamansampul
\halamanpengesahan
\paktaintegritas

\ringkasaneksekutif
Laporan ini disusun sebagai pemenuhan tugas mata kuliah Perancangan Sistem Teknik Fisika 1 (PSTF 1). Fokus utama proyek ini adalah perancangan sistem instrumentasi dan kontrol Proses ZXY pada Plant XYZ berdasarkan kebutuhan yang tertuang dalam dokumen lelang (TOR). Kelompok telah melakukan kajian mendalam mulai dari pemodelan matematis, desain rekayasa dasar, hingga analisis kelayakan ekonomi.

Hasil rancangan mencakup integrasi strategi kendali PID dan pengamanan proses. Identifikasi pasangan \textit{Manipulated Variable} (MV) dan \textit{Controlled Variable} (CV) telah dilakukan secara presisi untuk menekan \textit{Disturbance} (DV). Selain itu, kajian keselamatan menggunakan metode HAZOP memastikan risiko operasional berada dalam batas toleransi. Sistem yang dirancang dievaluasi secara finansial (CAPEX dan OPEX) guna menjamin nilai investasi yang kompetitif di tingkat industri.
\endringkasaneksekutif

\tableofcontents
\clearpage
\pagenumbering{arabic}

\chapter{KONSEPTUALISASI}

\section{Latar Belakang Dokumen Lelang TOR}
Dokumen \textit{Terms of Reference} (TOR) atau Kerangka Acuan Kerja merupakan instrumen krusial dalam siklus hidup proyek rekayasa sistem. Dokumen ini menjadi pedoman awal yang mendefinisikan batasan, tujuan, serta standar kualitas yang diharapkan oleh pemilik proyek kepada kontraktor atau perancang. Tanpa TOR yang jelas, perancangan sistem teknik fisik akan kehilangan arah dan berpotensi mengalami kegagalan fungsi saat diimplementasikan.

Dalam konteks Proses ZXY di Plant XYZ, dokumen lelang menekankan pada urgensi peningkatan efisiensi produksi sekaligus penguatan keandalan instrumen operasional. Hal ini dilatarbelakangi oleh kondisi instrumentasi pabrik yang sudah mulai usang serta perlunya adaptasi terhadap fluktuasi pasokan bahan baku. Oleh karena itu, tim perancang dituntut untuk memberikan solusi teknologi terkini yang adaptif dan presisi tinggi.

Permasalahan utama yang diidentifikasi dari TOR adalah ketidakstabilan variabel suhu dan tekanan pada saat proses perpindahan panas, yang seringkali memengaruhi kualitas produk akhir. Hal ini tidak hanya menurunkan efisiensi konversi massa, tetapi juga berimbas pada pemborosan energi. Kelompok merespons tantangan ini dengan mengusulkan perancangan ulang arsitektur kendali yang jauh lebih canggih dan responsif.

Selain aspek teknis kontrol, dokumen lelang juga menetapkan jadwal penyelesaian perancangan dan integrasi sistem yang cukup ketat. Hal ini menuntut manajemen proyek yang disiplin dari tim penyusun agar setiap fase (dari \textit{basic engineering} hingga \textit{commissioning}) berjalan selaras. Keterlambatan proyek akan berakibat pada penalti finansial yang signifikan bagi pihak integrator.

Lebih jauh lagi, TOR mensyaratkan ketersediaan sistem purna jual yang mudah diakses secara lokal. Keputusan ini dirancang untuk meminimalkan durasi perbaikan alat rusak (\textit{downtime}) pada operasi masa depan. Berdasarkan poin ini, kelompok telah memilih spesifikasi instrumen dengan ketersediaan suku cadang nasional dan dukungan vendor lokal terkemuka.

Pada akhirnya, dokumen TOR bertindak sebagai kerangka evaluasi bagi keseluruhan proyek ini. Bab-bab selanjutnya dalam laporan ini dirancang secara sistematis untuk menjawab seluruh tantangan teknis yang dimuat dalam latar belakang tersebut, memastikan sistem yang diserahkan memiliki rasio \textit{cost-to-benefit} yang maksimal.

\section{Penurunan Teori Model Sains Sistem}
Pemodelan sistem merupakan langkah pertama yang tak terelakkan dalam perancangan instrumentasi dan kontrol. Dengan membangun model matematis yang kuat, kita dapat memprediksi perilaku sistem terhadap berbagai gangguan (\textit{disturbances}) tanpa harus melakukan eksperimen \textit{trial-and-error} pada sistem fisik yang memakan biaya besar. Model ini diturunkan berlandaskan prinsip-prinsip konservasi massa dan energi dasar.

Kelompok kami mengawali penurunan dengan memformulasikan persamaan diferensial untuk dinamika aliran dan transfer panas pada Proses ZXY. Sebagai ilustrasi, hukum pertama termodinamika diaplikasikan untuk merumuskan profil perubahan suhu spesifik material terhadap waktu. Persamaan non-linear yang kompleks ini kemudian dilinearisasi menggunakan pendekatan deret Taylor di sekitar titik \textit{steady state}.

Transformasi Laplace kemudian dimanfaatkan untuk mengonversi persamaan diferensial kawasan waktu tersebut menjadi fungsi transfer dalam kawasan-$s$. Fungsi transfer inilah yang nantinya krusial dalam menganalisis perilaku respons transien, \textit{overshoot}, dan \textit{settling time} sistem di bawah berbagai mode kontrol (P, PI, atau PID). Evaluasi akar-akar pole juga dijalankan untuk verifikasi stabilitas mutlak sistem.

Selain murni berbasis persamaan analitis (\textit{first-principle modeling}), kami juga mengikutsertakan parameter-parameter empiris seperti efisiensi termal dan \textit{friction factor} perpipaan yang diakses dari spesifikasi alat aslinya. Kolaborasi teori fisika murni dengan konstanta empiris ini menjamin agar model matematis yang terbentuk tidak hanya indah secara teori, melainkan juga sangat relevan dengan realitas di pabrik.

Parameter dinamis yang diidentifikasi secara khusus mencakup \textit{Time Constant} ($\tau$) dan \textit{Dead Time} ($\theta$). Kehadiran \textit{dead time} dominan pada pipa transportasi cairan membutuhkan kompensasi kontrol khusus, karena algoritma PID konvensional umumnya tidak stabil untuk meredam osilasi akibat kelambatan sistem transportasi tersebut.

Bab penurunan teori ini merupakan bukti kokohnya pijakan keilmuan Teknik Fisika dalam proyek ini. Hasil pemodelan analitis inilah yang menjadi fondasi bagi pelaksanaan simulasi perangkat lunak pada tahapan desain selanjutnya, memastikan konfigurasi kontrol dapat mengatasi dinamika beban secara proaktif.

\section{Parameter Batas Operasional Teknis (URS)}
\textit{User Requirement Specifications} (URS) bertindak sebagai penerjemah dari keinginan operasional \textit{user} ke dalam parameter teknik kuantitatif yang mengikat. Parameter ini menjadi pembatas ketat dalam menentukan dimensi pipa, skala rentang instrumen pengukuran, hingga setpoint alarm di dalam program sistem terdistribusi (DCS). Tanpa URS, ruang lingkup desain dapat melebar di luar kebutuhan esensialnya.

Langkah pertama perumusan URS ini adalah mengkaji titik optimal operasional variabel fisika utama. Berdasarkan \textit{Process Flow Diagram} (PFD) pendahuluan, ditetapkan bahwa suhu dalam ruang pereaksi wajib dijaga pada $120^\circ\text{C}$ dengan deviasi absolut tidak melampaui $\pm 1.5^\circ\text{C}$. Rentang presisi operasional ini lantas digunakan sebagai acuan pemilihan resolusi pemancar suhu (\textit{temperature transmitter}).

Selain kondisi nominal operasional, batasan kegawatdaruratan juga tertuang secara definitif pada dokumen URS. Diwajibkan adanya mekanisme \textit{Emergency Shutdown System} (ESD) apabila batas tekanan melampaui 120\% dari tekanan desain, yakni lebih dari 6.5 bar. Batasan protektif ini tidak boleh dapat di-intervensi (\textit{override}) secara manual oleh pihak operator lapangan demi alasan keamanan.

Kondisi antarmuka atau \textit{Human Machine Interface} (HMI) pun tidak luput dari spesifikasi URS. Layar grafis diharuskan mampu merespons \textit{refresh rate} operasional di bawah 500 milidetik agar insiden keterlambatan observasi lapangan dapat ditekan. Tata letak layar dan palet warna alarm harus sepenuhnya mengadopsi prinsip-prinsip ergonomi dari panduan standar industri seperti ISA-101.

Dari segi ketahanan lingkungan, setiap \textit{field instrument} dituntut untuk patuh terhadap rating minimal IP-65 (Tahan Debu dan Semburan Air). Hal ini diputuskan mengingat kondisi lingkungan di Plant XYZ yang terpapar tingkat kelembaban tinggi secara kontinyu. Selain itu, perangkat harus tersertifikasi spesifik untuk zona area berbahaya (\textit{hazardous area classification}).

Sebagai rangkuman, URS adalah "konstitusi" perancangan. Kepastian angka dan parameter yang diikat secara definitif dalam bab ini mempermudah kelompok dalam membuat keputusan teknis terkait pengadaan instrumen pada tahapan \textit{Basic Engineering} dan pengujian pada saat \textit{Factory Acceptance Test}.

\chapter{FUNCTIONAL DESIGN}

\section{Logika Pergerakan Proses Kendali}
Logika pergerakan proses kendali adalah cetak biru "otak" sistem yang menggambarkan dengan terperinci bagaimana algoritma kontrol bereaksi atas rangsangan masukan dan gangguan lingkungan. Strategi pengendalian inti yang diterapkan pada Proses ZXY ini merupakan fusi fungsional antara \textit{feedback control} untuk meminimalkan deviasi lambat, dan strategi \textit{cascade control} untuk mengatasi fluktuasi parameter pendukung secara seketika.

Teks sekuensial yang disusun merujuk langsung pada standarisasi dokumen operasional ISA-5.06.01, yang diawali dari skenario penyalaan sistem (\textit{cold startup}). Dalam skenario ini, program memastikan kalibrasi dasar tervalidasi dan bahwa semua elemen katup perantara (\textit{control valves}) pada P\&ID berada dalam orientasi \textit{fail-safe} sebelum injeksi fluida material proses dilaksanakan.

Pusat arsitektur kontrol bergantung pada identifikasi presisi relasi antara \textit{Controlled Variable} (CV), \textit{Manipulated Variable} (MV), dan \textit{Disturbance Variable} (DV). Pada desain ini, logika komparasi error secara berkelanjutan dimonitor oleh blok kontrol internal DCS. Ketika DV seperti suhu lingkungan eksternal menyebabkan penurunan pada CV (misalnya suhu reaktor), kontroler akan seketika mengeksekusi perhitungan aksi integral-derivatif. Hasil perhitungan ini disalurkan dalam wujud keluaran sinyal 4-20 mA untuk menambah bukaan katup uap panas.

Tahapan logika paling tinggi ditutup oleh sistem \textit{interlock} regulatori. Dalam hal terjadinya deviasi anomali luar biasa—misalnya rusaknya sensor dan tekanan bejana menembus limit alarm atas—DCS akan menangguhkan \textit{loop} normal dan beralih paksa memutus katup suplai bahan. Skenario langkah-demi-langkah ini direkam dalam dokumen konfigurasi logika untuk diimplementasikan ke bahasa program berbasis \textit{Function Block Diagram} (FBD).

\chapter{BASIC ENGINEERING}

\section{Matriks Analitis Pembobot Opsi Desain}
Dalam ilmu rekayasa teknis, jarang sekali ditemukan sebuah skenario tunggal untuk memecahkan sebuah dilema desain operasional pabrik. Guna mencari keseimbangan strategis antara teknologi, keselamatan, dan investasi, tim menggunakan alat ukur objektif berupa Matriks Analitis Pembobot. Alat ini mereduksi preferensi personal desainer ke dalam formulasi rasional sistem skoring matematis yang adil.

Kami memulai evaluasi ini dengan merinci kriteria pembobot berdasarkan urgensinya bagi Plant XYZ. Kriteria 'Keandalan Instrumen (MTBF)' dan 'Kompatibilitas Interoperabilitas DCS' diberikan persentase pembobotan dominan, yakni masing-masing bernilai 30\%. Sementara itu, faktor 'Harga Modal Pembelian Awal' hanya diberi nilai 15\%, merujuk pada prinsip bahwa operasional jangka panjang lebih berharga ketimbang alat murah namun rawan kerusakan.

Matriks tersebut kemudian diaplikasikan untuk menyeleksi teknologi pengukur level material di reaktor: opsi A (Sensor Radar Ultrasonik) dibandingkan dengan opsi B (Sensor Tekanan Diferensial konvensional). Spesifikasi teknis vendor, catatan kinerja komersial, serta referensi literatur rekayasa digunakan sebagai dasar untuk memberikan poin mentah skala 1 hingga 5 pada setiap kategori pengujian.

Melalui perkalian skoring dengan bobot parameternya, Opsi A (Radar) terpilih dengan kemenangan marjinal meskipun biaya investasinya lebih tinggi. Penilaian memenangkan sensor radar dikarenakan faktor imunitasnya terhadap fluktuasi densitas cairan yang sering bergejolak, yang merupakan kelemahan mutlak opsi B. Proses analisis matriks ini menjadi justifikasi valid perancangan yang kuat secara argumentasi rekayasa di mata direksi pabrik.

Dokumentasi dari skoring analitik ini akan diarsipkan sebagai lampiran resmi perhitungan proyek. Rekam jejak transparansi evaluasi desain opsi awal memastikan bahwa perancangan Proses ZXY bukan berlandaskan pada metodologi acak, namun melalui kerangka pemikiran rekayasa komparatif. Keputusan yang didapat dari langkah ini secara sah menjadi spesifikasi instrumen definitif untuk pesanan komponen.

\section{Process Flow Diagram (PFD)}
Representasi arsitektur mekanikal dan kimiawi paling mendasar bagi Plant XYZ adalah \textit{Process Flow Diagram} (PFD). Dokumen visual ini memaparkan tata urutan unit-unit alat berat—mulai dari bejana penampung awal, reaktor reaksi bersuhu tinggi, kolom distilasi pemisah, hingga pendingin produk akhir. PFD adalah landasan orientasi komprehensif tanpa harus terjebak ke dalam rumitnya simpangan pipa mikro dan katup-katup isolasi minor.

Seluruh simbologi bentuk alat yang dicitrakan telah dipastikan mematuhi kodifikasi standar ISO-10628. Penggunaan lambang internasional ini menghilangkan risiko bias persepsi antara teknisi sipil, insinyur proses, maupun spesialis teknik kendali. Selain peralatan sentral, jalur utilitas pendukung (seperti udara instrumen tekan, pasokan nitrogen murni, dan utilitas air pendingin) juga digambarkan secara makro.

Satu aspek tak terpisahkan dari PFD adalah penyertaan indikasi titik kontrol kunci pada jalur fluida massal. Hal ini dimanfaatkan oleh tim teknik fisika untuk mulai mendefinisikan lokasi spasial yang paling representatif untuk peletakan sensor instrumen. Sebagai contoh, letak elemen pengukur tekanan tidak boleh terganggu oleh turbulensi lokal di dekat sambungan pompa sirkulasi.

Lebih lanjut, dalam PFD dipresentasikan penanda arah panah aliran beserta nomor identifikasi percabangan (\textit{stream tags}). Tag nomor aliran inilah yang mengikat secara logis gambar visual PFD dengan tabulasi tabel neraca massa kuantitatif pada subbab selanjutnya. Visualisasi relasional ini menyempurnakan pemahaman perihal transformasi fisis dari bahan baku mentah menjadi produk komersial jadi.

\section{Neraca Massa dan Neraca Termal}
Tanpa landasan data perpindahan massa dan distribusi panas yang mumpuni, sistem pengukuran kontrol hanyalah infrastruktur buta yang beroperasi tanpa parameter batas desain yang valid. Oleh karena itu, tabulasi teoritis dari Neraca Massa dan Neraca Energi (\textit{Heat Balance}) memegang peranan mutlak bagi tahapan penskalaan presisi katup kendali (\textit{control valve sizing}) serta verifikasi kapasitas operasional proses.

Dalam penyusunan neraca massa, asas universal $\sum M_{in} = \sum M_{out} + \Delta M_{acc}$ dipegang teguh. Sebagai misal perhitungan teoritis pada sebuah subsistem percampuran, apabila aliran Umpan A terhitung mensuplai 500 kg/jam dan Umpan B menyuntikkan 300 kg/jam, maka laju alir keluaran produk ideal akan dijaga konstan di titik 800 kg/jam pada rezim tunak (tanpa kerugian bocoran).

Bersamaan dengan itu, tinjauan konversi suhu dilakukan menggunakan kalkulasi Termodinamika Entalpi ($H$). Dengan mensubstitusikan kalor spesifik cairan ($C_p$), suhu masukan yang bervariasi dari tiap cabang pipa dipetakan sehingga estimasi profil suhu ekuilibrium keluaran dapat diprediksi secara empirik. Ringkasan tabulasi komprehensif antara sifat massa dan energi reaktor tersaji pada Tabel \ref{tab:neracamassa}.

\begin{table}[htbp]
\centering
\caption{Ringkasan Neraca Massa dan Energi pada Reaktor ZXY}
\label{tab:neracamassa}
\resizebox{\textwidth}{!}{%
\begin{tabular}{|p{5cm}|c|c|c|c|}
\hline
\textbf{Identitas Aliran (Stream)} & \textbf{Suhu ($^\circ$C)} & \textbf{Tekanan (bar)} & \textbf{Massa (kg/jam)} & \textbf{Entalpi (kJ/jam)} \\ \hline
Masuk 1 (Bahan Baku Utama) & 25.0 & 1.01 & 500.0 & 12,500 \\ \hline
Masuk 2 (Aditif Kimiawi) & 30.0 & 1.01 & 300.0 & 9,000 \\ \hline
\textbf{Keluaran (Produk Reaktor)} & \textbf{85.0} & \textbf{1.01} & \textbf{800.0} & \textbf{35,000} \\ \hline
Air Pendingin Utilitas (Inlet) & 25.0 & 2.00 & 1,200.0 & - \\ \hline
Air Pendingin Utilitas (Outlet) & 45.0 & 1.80 & 1,200.0 & - \\ \hline
\end{tabular}%
}
\end{table}

Data aliran volumetrik dari Tabel \ref{tab:neracamassa} merupakan fondasi bagi penentuan parameter kapasitas katup ($C_v$ atau \textit{Flow Coefficient}). Katup-katup pada plant akan dirancang kemampuannya agar optimal pada skenario aliran normal (800 kg/jam) namun tetap sanggup beroperasi secara linear apabila laju produksi di-eskalasi mencapai 120\% dari batas wajar saat diperlukan (\textit{maximum rated flow}).

Kalkulasi entalpi pada sistem pendingin menegaskan besaran kalor buangan reaksi yang harus secara konstan dibuang dari lingkungan reaktor. Data beda suhu pendingin inilah yang menuntun kami pada spesifikasi resolusi sensor \textit{RTD (Resistance Temperature Detector)} yang direkomendasikan untuk loop pendinginan proses.

\section{Identifikasi Variabel Kendali}
Berdasarkan parameter operasional neraca massa dan energi yang telah dihitung sebelumnya, ketepatan pemasangan antar variabel kendali wajib didefinisikan agar proses tetap reaktif terhadap perturbasi eksternal dan tidak memasuki siklus instabilitas. Identifikasi relasi presisi antara \textit{Controlled Variable} (CV), \textit{Manipulated Variable} (MV), dan \textit{Disturbance Variable} (DV) dapat dilihat secara mendetail pada Tabel \ref{tab:mvcV}.

\begin{table}[htbp]
\centering
\caption{Identifikasi Variabel Kendali (CV, MV, DV) pada Proses ZXY}
\label{tab:mvcV}
\resizebox{\textwidth}{!}{%
\begin{tabular}{|c|p{4.5cm}|p{5cm}|p{4.5cm}|}
\hline
\textbf{No. Loop} & \textbf{Controlled Variable (CV)} & \textbf{Manipulated Variable (MV)} & \textbf{Disturbance Variable (DV)} \\ \hline
TC-101 & Suhu Utama Reaktor ($^\circ$C) & Laju Alir Uap Pemanas (kg/s) melalui TV-101 & Perubahan Suhu Umpan, Suhu Lingkungan \\ \hline
PC-102 & Tekanan \textit{Headspace} (bar) & Bukaan Katup Ventilasi / PV-102 (\%) & Lonjakan Reaksi Kimia Internal \\ \hline
LC-103 & Level Cairan Pemisah (m) & Laju Aliran Produk Bawah melalui LV-103 & Fluktuasi Kapasitas Laju Umpan Utama \\ \hline
FC-104 & Rasio Umpan Kimia & Putaran Pompa Dosis / VFD-104 (RPM) & Fluktuasi Hambatan/Viskositas Aliran \\ \hline
\end{tabular}%
}
\end{table}

Matriks relasi pada Tabel \ref{tab:mvcV} merupakan kerangka dasar dalam perancangan strategi kontrol \textit{Multi-Input Multi-Output} (MIMO) apabila terjadi peleburan interaksi antar loop. Tabel ini juga menjadi pijakan dalam menentukan prioritas pengendalian (misal: pengaturan tekanan harus lebih cepat dari suhu).

\section{Penjelasan Alur Draf P\&ID}
\textit{Piping and Instrumentation Diagram} (P\&ID) mentranslasikan abstrak desain PFD menjadi realitas kerangka teknis perpipaan, pengkabelan fungsional, serta posisi absolut dari seluruh pemancar deteksi (\textit{sensors}). Dalam lembar kerja ini (yang akan direproduksi dalam ukuran kertas A3), hierarki pergerakan fluida dan skema aksi sinyal kontrol dapat dibaca layaknya cetak biru anatomis dari pembuluh nadi sebuah pabrik industri canggih.

Metodologi penamaan yang disisipkan sangat tunduk pada notasi baku The International Society of Automation (ISA-5.1). Hal ini dibuktikan melalui pemberian kodifikasi bola (*bubble tags*) lengkap. Contohnya, "FT-201" bermakna bahwa posisi fisik perangkat adalah instrumen pengukur debit luapan di area unit 200, sedangkan "FV-201" menerangkan posisi elemen penutup akhir katup mekanik yang menerima sinyal operasinya.

Garis penghubung antar balok kontrol dilukiskan dengan representasi yang memiliki filosofi pembeda—garis solid ganda untuk jalur transportasi pneumatik bertenaga angin, sedangkan garis putus-putus bergelombang merepresentasikan transfer data kelistrikan 4-20 mA analog maupun komunikasi bus digital. Konsistensi penarikan garis interkoneksi ini membebaskan instalatur lapangan dari miskonsepsi orientasi pengkabelan \textit{hardwiring}.

Titik sentral kejelasan draf ini terletak pada ilustrasi aksi keadaan gagal instrumentasi mekanik (\textit{Failure Mode}). Saat kehilangan tegangan listrik atau sumber suplai kompresor pneumatik, simbol katup diberi inisial FC (\textit{Fail Close}) atau FO (\textit{Fail Open}). Strategi mode kegagalan ini selaras dengan profil keselamatan lingkungan yang menjamin proses bergeser mengamankan diri (menutup suplai reagen berbahaya otomatis) bila pasokan daya plant lenyap.

Di ujung perancangan sistem, draf P\&ID diposisikan sebagai jembatan negosiasi terakhir antardepartemen sebelum finalisasi \textit{Approved for Construction} (AFC). Setiap revisi yang dibahas secara multidisiplin selama perumusan draf disepakati melalui persetujuan paraf. Hal tersebut menunjukkan peran substansial \textit{P\&ID} bukan saja sekadar dokumen pelengkap, melainkan nyawa instruksional utama perancangan sistem.

\chapter{KAJIAN KESELAMATAN}

\section{Hazard and Operability Study (HAZOP)}
Risiko adalah inherent di setiap pabrik dengan pengolahan fisika dan kimia dinamis; membiarkannya tanpa analisis terstruktur setara dengan menyusun skenario bencana di masa depan. Kelompok perancang mengutilisasi teknik deduktif komprehensif berupa Metode \textit{Hazard and Operability Study} (HAZOP) demi membongkar titik lemah sistem sebelum diwujudkan ke dalam wujud fisik. HAZOP menggunakan stimulasi linguistik yang dikenal sebagai panduan kata (*Guide Words*).

Mekanisme audit dilaksanakan dengan mensegmentasi plant menjadi area analisis mikro atau "Simpul" (Nodes). Tiap komponen kunci di setiap Node dihadapkan dengan kata sandi analitis seperti "NO", "LESS", "MORE", "REVERSE", dan seterusnya. Jika disandingkan pada parameter dasar, misalnya laju massa umpan reaktor, kata "LESS" akan menstimulasi pertanyaan: \textit{Apa yang terjadi jika pasokan fluida tiba-tiba berkurang tajam?}

Proses evaluasi sebab-akibat kemudian diidentifikasi dan ditabulasi. Kegagalan dari sebuah skenario lalu ditakar dampaknya terhadap stabilitas reaksi dan potensi cedera personal operator atau kontaminasi lingkungan, merujuk pada standar matriks risiko. Visualisasi tabulasi dari salah satu penelaahan titik kritis tersebut dihadirkan pada Tabel \ref{tab:hazop}.

\begin{table}[htbp]
\centering
\caption{Contoh Lembar Kerja HAZOP (Node: Titik Pipa Umpan Utama Reaktor)}
\label{tab:hazop}
\resizebox{\textwidth}{!}{%
\begin{tabular}{|l|l|p{4.5cm}|p{4.5cm}|p{4.5cm}|}
\hline
\textbf{Parameter} & \textbf{Guide Word} & \textbf{Penyebab Kegagalan} & \textbf{Konsekuensi Lanjutan} & \textbf{Aksi Proteksi / Rekomendasi} \\ \hline
\textbf{Aliran} & LESS / NO & Katup tersumbat endapan, pompa tiba-tiba mati. & Rasio terganggu, pendinginan lemah, proses terhenti. & Alarm aliran lambat (FAL), Interlock menghidupkan pompa cadangan. \\ \hline
\textbf{Suhu} & MORE & Gangguan suplai \textit{Cooling Water} hulu. & Reaktor \textit{overheating}, potensi eksotermik liar (\textit{runaway}). & Implementasi Sensor Redundan (TT), \textit{High-High Temp Alarm} (TAHH). \\ \hline
\textbf{Tekanan} & HIGHER & Fasa gas berlebih terakumulasi di bejana. & Tegangan batas material tangki melampaui limit, risiko pecah. & PSV dipasang secara mekanikal langsung melepas tekanan ke sirkuit \textit{Flare}. \\ \hline
\end{tabular}%
}
\end{table}

Berdasarkan Tabel \ref{tab:hazop}, terlihat secara persis rekomendasi instalasi protektif apa saja yang dilahirkan akibat identifikasi bahaya ini. Seluruh usulan perlindungan instrumen mandiri tambahan dari kertas kerja HAZOP lalu langsung disirkulasikan sebagai perintah modifikasi ulang (*revision cycle*) pada lembar draf akhir P\&ID yang sudah disinggung sebelumnya. 

Tanggung jawab moral sekaligus profesional insinyur teknik fisika diletakkan utuh pada fase ini. Kajian HAZOP membuktikan secara empiris bahwa kelompok tidak hanya merancang proses kendali untuk kondisi cerah benderang (*sunny day conditions*), namun secara radikal juga merancang arsitektur keamanan tingkat lanjut demi meredam kerugian operasional dan menjaga nyawa pekerja pabrik.

\section{Penentuan Safety Integrity Level (SIL)}
Melengkapi identifikasi kualitatif HAZOP, perhitungan deterministik berskala numerik dihadirkan melalui penentuan ambang batas \textit{Safety Integrity Level} (SIL). Tingkatan kuantifikasi SIL ditetapkan berdasarkan regulasi payung arsitektur keselamatan fungsional IEC 61508 serta pedoman turunan industri proses IEC 61511. Level integrasi ini mengukur sejauh mana fungsional perlindungan instrumentasi sanggup menekan statistik angka celaka probabilistik tahunan.

Melalui penerapan model analitis grafik risiko (*Risk Graph Method*), parameter fatalitas korban jiwa, lamanya paparan radiasi/bahaya, serta kemungkinan evakuasi dini turut dinilai secara objektif. Tinjauan mendalam membuktikan bahwa *Safety Instrumented Function* (SIF) khusus pencegahan tekanan ledak pada area Reaktor Proses ZXY dituntut untuk minimal memiliki level kinerja keselamatan SIL 2.

Meraih standar spesifikasi SIL 2 bukanlah persoalan konfigurasi parameter secara digital, melainkan menuntut reformasi perangkat keras pengindra lapangan. Kelompok kami menugaskan instalasi sensor \textit{redundant} ganda dengan metode seleksi pemungutan perbandingan suara algoritma "1oo2" (*One Out Of Two*). Bila satu sensor membisu atau error, sistem pengamanan tetap diizinkan bereaksi melalui kepastian operasional dari sisa sensor hidup.

Angka kemungkinan sistem lumpuh pada saat situasi kritis tengah terjadi atau dikenal sebagai besaran *Probability of Failure on Demand* ($PFD_{avg}$) diverifikasi secara analitis dengan formula reliabilitas Markov. Rasio kerusakan suku cadang sensor, kendali \textit{Logic Solver}, serta aktuator pneumatik akhir ditotal keseluruhannya dan dijaga ketat agar selalu di bawah batas toleransi deviasi probabilistik numerik SIL 2.

Tahapan rekayasa berkesinambungan ini juga merekomendasikan perlunya pengujian lapangan repetitif per siklus 12 bulan (\textit{Proof Test Interval}). Tes validasi tahunan dipetakan dalam prosedur baku guna merawat angka reliabilitas $PFD_{avg}$ instrumen agar grafiknya tidak merosot anjlok. Keteguhan implementasi SIL ini merupakan perlambang integrasi manajemen pemeliharaan proaktif terhadap komitmen keselamatan yang bernilai jual tinggi.

\chapter{KAJIAN TEKNO-EKONOMI}

\section{Capital Expenditure (CAPEX)}
Tingkat kecanggihan sistem instrumentasi fisika wajib diimbangi oleh realitas ketersediaan kas injeksi tunai dan kapabilitas persetujuan modal awal atau disebut \textit{Capital Expenditure} (CAPEX). Perhitungan modal di awal ini menjadi dokumen justifikasi kelayakan utama bagi pihak manajemen perusahaan atau dewan direktur sebelum mereka membubuhi persetujuan eksekusi konstruksi (\textit{Final Investment Decision}).

Pembentukan lembar anggaran CAPEX ini menelusuri riwayat seluruh peralatan dalam \textit{Bill of Materials} (BoM) yang diproduksi dari revisi final spesifikasi dasar di Bab 3. Pendekatan valuasi biaya diperoleh dari kueri referensi pasar, riwayat proyek pabrik percontohan yang setara, dan estimasi pabrikasi komersial. Tabulasi rincian prakiraan penyerapan dana ditunjukkan secara transparan pada Tabel \ref{tab:capex}.

\begin{table}[htbp]
\centering
\caption{Rincian Estimasi Perolehan Dana Capital Expenditure (CAPEX)}
\label{tab:capex}
\resizebox{\textwidth}{!}{%
\begin{tabular}{|c|l|c|r|}
\hline
\textbf{No} & \multicolumn{1}{c|}{\textbf{Klasifikasi Deskripsi Item Utama}} & \textbf{Kuantitas} & \textbf{Total Biaya (USD)} \\ \hline
1 & Kontroler Utama (PLC/DCS) dengan modul Redundansi & 1 Set Panel & 45,000 \\ \hline
2 & Sistem Sensor Instrumen Pengindraan Berstandar SIL 2 & 12 Unit & 18,500 \\ \hline
3 & Elemen Aktuator Katup Pneumatik (\textit{Control Valves}) & 8 Unit & 24,000 \\ \hline
4 & Kabel Sinyal Industri, Jalur Tray, dan Komponen Penunjang & 1 Paket Lot & 8,500 \\ \hline
5 & Alokasi Jasa Tenaga Engineering, Uji Integrasi, dan Instalasi & 1 Paket Jasa & 20,000 \\ \hline
\multicolumn{3}{|r|}{\textbf{Akumulasi Total Estimasi CAPEX Proyek (Awal)}} & \textbf{116,000} \\ \hline
\end{tabular}%
}
\end{table}

Berdasarkan kompilasi angka Tabel \ref{tab:capex}, pembobotan kapitalisasi investasi perangkat lunak (DCS) memang cukup menggerus separuh porsi belanja, hal ini diakibatkan oleh pemilihan arsitektur otak redundansi demi alasan keandalan total tanpa henti (\textit{zero fault tolerance}). Alokasi instalasi mencakup segala ongkos sertifikasi kalibrasi pengawasan \textit{Site Acceptance Test} (SAT) oleh tenaga teknik lapangan.

Kami menambahkan injeksi dana kontinjensi sebesar 12\% sebagai tameng risiko komersial. Cadangan ini krusial untuk menyerap guncangan biaya tersembunyi yang lazim terjadi seperti fluktuasi indeks kurs mata uang pengadaan, kemacetan kargo rantai pasok maritim global instrumen, hingga taksiran kenaikan upah perakitan mendadak. CAPEX adalah kerangka arsitektur disiplin keuangan demi mewujudkan gagasan ilmiah kendali perancangan teknis.

\section{Operational Expenditure (OPEX)}
Tidak ada jaminan stabilitas keuntungan jangka panjang pada seumur hidup proses plant pabrik tanpa perhitungan alokasi operasional perawatan atau \textit{Operational Expenditure} (OPEX) yang rasional. OPEX merupakan sirkulasi beban tunai teratur berkesinambungan agar infrastruktur automasi, kontrol kendali, perangkat keras kelistrikan, maupun daya pikir pekerja teknis lapangan tetap melaju di batas optimal kapasitas kerjanya.

Berbeda dengan proyek statis bangunan sipil, jantung kelangsungan sistem kendali berpusat pada ketersediaan listrik dan konsumsi kompresor udara tekanan instrumen untuk menghidupkan fungsi gerak elemen akhir tiap katup. Ongkos dari kebutuhan fluida daya dorong yang menyala non-stop 24 jam sehari selama 330 hari kerja ini dikuantifikasi dalam perhitungan daya tahunan. Selain listrik, rincian proyeksi biaya OPEX disarikan pada Tabel \ref{tab:opex}.

\begin{table}[htbp]
\centering
\caption{Proyeksi Estimasi Operational Expenditure (OPEX) Tahunan Plant XYZ}
\label{tab:opex}
\resizebox{\textwidth}{!}{%
\begin{tabular}{|c|l|r|}
\hline
\textbf{No} & \multicolumn{1}{c|}{\textbf{Deskripsi Kategori Biaya}} & \multicolumn{1}{c|}{\textbf{Proyeksi Tahunan (USD)}} \\ \hline
1 & Konsumsi Pemakaian Energi Listrik dan Udara Tekan Harian & 8,000 \\ \hline
2 & Eksekusi Pemeliharaan Preventif dan Kalibrasi Sensor Independen & 5,500 \\ \hline
3 & Manajemen Pergantian Komponen Suku Cadang (\textit{Consumables}) & 3,000 \\ \hline
4 & Biaya Pembaruan Lisensi Perangkat Lunak dan Bantuan \textit{IT Support} & 2,500 \\ \hline
5 & Pelatihan Sertifikasi Peningkatan Kemahiran Operator dan Teknisi & 2,000 \\ \hline
\multicolumn{2}{|r|}{\textbf{Total Estimasi OPEX per Siklus Tahun Operasional}} & \textbf{21,000} \\ \hline
\end{tabular}%
}
\end{table}

Kebutuhan penggantian \textit{consumable parts} diurutkan pada siklus pakai dari vendor (Tabel \ref{tab:opex} Item 3). Suku cadang \textit{O-rings}, pelumas stem katup mekanik mekanis, dan filter udara diganti presisi secara terjadwal meskipun secara mata fisik instrumen tersebut terkesan wajar. Pergantian preventif merupakan metodologi OPEX efisien demi membunuh probabilitas rusaknya unit mesin utama senilai puluhan ribu dolar hanya gara-gara mengirit suku cadang murah senilai beberapa dolar.

Komitmen pemeliharaan kompetensi operator turut difasilitasi dalam lini pendanaan OPEX rutin. Tidak peduli seberapa revolusionernya DCS dan P\&ID yang didirikan oleh tim desainer instrumentasi teknis, keseluruhan sistem tetap rawan hancur karena malapraktik tenaga pekerja kaku di ruang panel kontrol. Mengalokasikan dana untuk pelatihan berkelanjutan operator lokal mencerminkan visi kematangan desainer menyongsong teknologi proses modern 5.0.

\section{Analisis Kelayakan Finansial}
Arsitek sejati perancangan teknik fisika adalah mereka yang mengolaborasikan intuisi rekayasa tingkat tinggi di hadapan realitas kalkulus ekonomi investasi kelayakan. Di penghujung proposal rancangan, analisis penutup dilanjutkan dengan menakar kriteria \textit{Net Present Value} (NPV) secara matematis yang bertujuan menyandingkan ekuivalen dana cair yang diinvestasikan masa kini dengan keuntungan tahunan penghematan di dekade operasi masa depan.

Melalui penerapan kalkulasi NPV, kami menetapkan nilai suku bunga persentase diskonto (\textit{Discount Rate}) konservatif sebesar 10\% dengan bentang horizon perhitungan fungsional selama masa pakai instrumen yakni kira-kira mencapai waktu produktif 10 tahun penuh. Evaluasi optimis dari proyek menunjukan angka perolehan akhir simulasi arus kas NPV yang berlabuh teguh di ranah nilai teritorial mata uang yang positif. 

Sebagai indikator pelengkap pertimbangan komite perusahaan, proyeksi kurun waktu \textit{Payback Period} diderivasikan. Kalkulasi ini bertugas melacak pada siklus bulan keberapa selisih efisiensi proses akibat instrumentasi baru kami ini mampu melunasi lunas tagihan besar CAPEX murni pada awalnya. Rekayasa automasi efisien Proses ZXY mencatatkan periode pelunasan titik impas (BEP) mengagumkan yakni pada triwulan di kurun durasi pemakaian operasi empat tahun pertama.

Analisis sensitivitas yang ketat terhadap deviasi indeks investasi juga dilakukan. Kami merekayasa kalkulasi pesimis atas skenario andaikata kenaikan OPEX meledak tidak terkendali hingga membebani kenaikan inflasi melebihi batas 15\%. Kendati ditempa asumsi fiskal pelemahan ekonomi tersebut, parameter \textit{Internal Rate of Return} (IRR) finansial instrumen kontrol tetap tegar berada aman di atas indeks suku bunga permodalan pokok pinjaman margin bank nasional rata-rata komersial korporasi reguler.

Kesimpulan akhir yang dihasilkan dalam analisis purna rancangan perancangan tekno-ekonomi adalah rekomendasi bulat yang kuat kepada investor jajaran staf ahli pabrik. Perancangan Sistem Instrumentasi dan Kontrol ZXY pada Plant XYZ tidak semata-mata unggul perihal spesifikasi penyelarasan keamanan dan kalibrasi teknologi \textit{advanced control logic}, melainkan ini jelas adalah skema \textit{business-case} yang sangat cemerlang dengan kelayakan profitabilitas masa depan cerah, reliabel, mandiri dan menjanjikan tingkat pengembalian keuntungan rasional operasional jangka panjang yang sangat memadai.

\appendix
\chapter{DRAF P\&ID}
Cetak biru gambaran skematik proses (Piping \& Instrumentation Diagram) serta tatanan tata letak penempatan instrumen-instrumen ukur Plant ZXY dapat dilihat dan dilampirkan secara utuh pada halaman ekstensi berskala besar kertas ukuran teknis (A3) di bagian ini.

\chapter{KUMPULAN SALINAN LOGBOOK}
Sebagai rekam jejak historis pemenuhan capaian tanggung jawab tugas mata kuliah Perancangan Sistem Teknik Fisika 1 (PSTF 1), salinan fotokopi otentik progres mingguan \textit{Logbook} riset lapangan dengan bubuhan stempel basah tanda tangan pembimbing pengesah dilampirkan pada bab verifikasi orisinalitas riset akhir kelompok ini.

\end{document}
